% Generated draft: requires mathematical and source review.
\documentclass[11pt]{article}
\usepackage[margin=25mm]{geometry}
\usepackage{fontspec}
\setmainfont{Latin Modern Roman}
\usepackage{amsmath,amssymb,mathtools,unicode-math}
\setmathfont{Latin Modern Math}
\usepackage{enumitem,needspace,xurl,hyperref}
\setlength{\parindent}{0pt}
\setlength{\parskip}{5pt}
\setlength{\emergencystretch}{3em}
\newcommand{\N}{\mathbb{N}}
\newcommand{\Z}{\mathbb{Z}}
\newcommand{\Q}{\mathbb{Q}}
\newcommand{\R}{\mathbb{R}}
\newcommand{\C}{\mathbb{C}}
\newcommand{\F}{\mathbb{F}}
\newcommand{\A}{\mathbb{A}}
\newcommand{\PP}{\mathbb{P}}
\newcommand{\E}{\mathbb{E}}
\newcommand{\eps}{\varepsilon}
\newcommand{\dd}{\,\mathrm{d}}
\begin{document}
\section*{UnsolvedMath ID: 701437}
\textbf{Source problem code:} catalogue entry
\subsection{Definitions and mathematical statement}
All groups under consideration are finite. A formation \(\mathfrak F\) is a nonempty class of finite groups closed under homomorphic images and finite subdirect products. A subdirect product of groups \(G_1,\ldots,G_n\) is a subgroup of \(\prod_{i=1}^n G_i\) whose projection onto each \(G_i\) is surjective. For a finite group \(H\), let \(\Phi(H)\) denote its Frattini subgroup, the intersection of its maximal proper subgroups, with \(\Phi(1)=1\).

Let \(\mathcal X\) be a class of finite simple groups. Write \(O_{\mathcal X}(G)\) for the largest normal subgroup of \(G\) all of whose composition factors belong to \(\mathcal X\). Here composition factors are the simple factors in a composition series, considered up to isomorphism. The trivial subgroup satisfies this condition vacuously.

The source imposes the hypothesis \(\pi(\mathcal X)=\operatorname{char}(\mathcal X)\) without defining either symbol. Under their conventional interpretation in this setting, these are
\[
\pi(\mathcal X)=\{p\text{ prime}:p\mid |S|\text{ for some }S\in\mathcal X\},\qquad
\operatorname{char}(\mathcal X)=\{p\text{ prime}:C_p\in\mathcal X\},
\]
where \(C_p\) is the cyclic group of order \(p\). This interpretation requires confirmation against the cited source.

For a class \(\mathcal X\) satisfying the stated equality, a formation \(\mathfrak F\) is \(\mathcal X\)-saturated if
\[
\text{for every finite group }G,\qquad
G/\Phi\bigl(O_{\mathcal X}(G)\bigr)\in\mathfrak F
\ \Longrightarrow\ G\in\mathfrak F.
\]
The question is whether, for every such class \(\mathcal X\) and every formation \(\mathfrak F\),
\[
\mathfrak F\text{ is }\mathcal X\text{-saturated}
\quad\Longrightarrow\quad
\mathfrak F\text{ is }\mathcal X\text{-local in the sense of F\"orster}.
\]
The supplied record refers to F\"orster's 1985 publication for the definition of \(\mathcal X\)-locality but does not reproduce it. Consequently, the conclusion is retained as a source-dependent predicate; a self-contained definition of that predicate remains unresolved.
\subsection{Short English statement}
For every class \(\mathcal X\) of finite simple groups satisfying the source's condition \(\pi(\mathcal X)=\operatorname{char}(\mathcal X)\), must every \(\mathcal X\)-saturated formation be \(\mathcal X\)-local in F\"orster's sense? The locality definition is absent from the supplied record, and the interpretation of the prime-set notation requires confirmation.
\subsection{Sources}
\begin{itemize}
\item Local UnsolvedMath registry, supplied record with permanent ID \(701437\), exact source code \texttt{catalogue entry}, title \emph{Kourovka Notebook Problem 18.8}. The supplied status is \texttt{open}; this label is not an independently verified mathematical assessment.
\item A. Ballester-Bolinches, proposer, 2014. The supplied background identifies Problem 18.8 in \emph{The Kourovka Notebook (Unsolved Problems in Group Theory)}, 18th issue (2014), and cites the 21st issue edition (2026), arXiv identifier \texttt{1401.0300v48}. The registry describes the collection as edited by E. I. Khukhro and V. D. Mazurov.
\item P. F\"orster, \emph{Publ. Sec. Mat. Univ. Aut\`onoma Barcelona}, \textbf{29}, no. 2--3 (1985), 39--76. The supplied statement directs the reader to this reference for the definition of \(\mathcal X\)-locality. No article title or URL is supplied, and the reference has not been independently consulted.
\end{itemize}
\end{document}
